\section{Physical \& Mathematical Framework}

\subsection{1D Hydrostatic Interior Structure}
Assuming spherical symmetry and instantaneous convective mixing along an envelope adiabat $S_{\mathrm{env}}$:
\begin{equation}
\frac{dr}{dm} = \frac{1}{4\pi r^2 \rho}, \quad \frac{dP}{dm} = -\frac{G m}{4\pi r^4}, \quad \frac{dT}{dm} = -\frac{G m T \nabla_{\mathrm{ad}}}{4\pi r^4 P}
\end{equation}

\subsection{Equations of State (EOS)}
\begin{itemize}
    \item \textbf{Central Core ($0 \le m \le M_c$)}: Solved via a 3rd-order Birch-Murnaghan isothermal EOS for high-pressure rock/ice mixtures.
    \item \textbf{Envelope ($M_c < m \le M_p$)}: Modeled via SCVH/CMS tabular data or non-ideal degenerate gas law.
\end{itemize}

\subsection{Atmospheric Boundary Conditions}
Coupled to Guillot (2010) semi-analytical irradiated atmosphere profile:
\begin{equation}
T^4(\tau) = \frac{3 T_{\mathrm{int}}^4}{4}\left(\tau + \frac{2}{3}\right) + \frac{3 T_{\mathrm{irr}}^4}{4}\left[\frac{2}{3} + \frac{2}{3\gamma}\left(1 + \left(\frac{\gamma \tau}{2} - 1\right)e^{-\gamma \tau}\right)\right]
\end{equation}

\subsection{Core Mass \& Metallicity Dependence}
Heavy-element core mass $M_c$ scales with host star metallicity $[\mathrm{Fe/H}]$ (Thorngren et al. 2016):
\begin{equation}
M_c([\mathrm{Fe/H}], M_p) \approx 15 \left(\frac{M_p}{M_{\mathrm{Jup}}}\right)^{0.6} 10^{0.5 [\mathrm{Fe/H}]} M_\oplus
\end{equation}

\subsection{Interior Energy Injection (Tidal \& Ohmic)}
\begin{align}
P_{\mathrm{tidal}} &= \frac{21}{2} \left(\frac{k_2}{Q}\right) \frac{G M_*^2 R_p^5 e^2}{a^6} \\
P_{\mathrm{Ohmic}} &= \epsilon_0 \exp\left[ -\frac{(T_{\mathrm{eq}} - 1600\text{ K})^2}{2 (300\text{ K})^2} \right] \pi R_p^2 F_{\mathrm{inc}} (1 - A_b)
\end{align}

\subsection{Coupled Orbital Element \& Spin Vector Evolution}
The 6D coupled dynamical-thermal state vector $\mathbf{y}(t) = [S_{\mathrm{env}}, a, e, I, \Omega_{\mathrm{rot}}, \varepsilon]^T$ evolves under equilibrium tidal dissipation (Hut 1981, Eggleton et al. 1998, Leconte et al. 2010):
\begin{align}
\frac{da}{dt} &= -\frac{57}{2} \left(\frac{k_2}{Q}\right) \frac{M_*}{M_p} \left(\frac{R_p}{a}\right)^5 n a e^2 + 3 \left(\frac{k_2}{Q}\right) \frac{M_*}{M_p} \left(\frac{R_p}{a}\right)^5 n a \left( \frac{\Omega_{\mathrm{rot}}}{n} \cos\varepsilon - 1 \right) \\
\frac{de}{dt} &= -\frac{21}{2} \left(\frac{k_2}{Q}\right) \frac{M_*}{M_p} \left(\frac{R_p}{a}\right)^5 n e \left[ 1 - \frac{11}{18} \frac{\Omega_{\mathrm{rot}}}{n} \cos\varepsilon \right] \\
\frac{d\Omega_{\mathrm{rot}}}{dt} &= -\frac{3}{2 C_p M_p R_p^2} \left(\frac{k_2}{Q}\right) \frac{G M_*^2 R_p^5}{a^6} \left(\Omega_{\text{rot}} - n \cos\varepsilon\right) - \frac{2 \Omega_{\mathrm{rot}}}{R_p} \frac{dR_p}{dt} \\
\frac{d\varepsilon}{dt} &= -\frac{3}{2 C_p M_p R_p^2} \left(\frac{k_2}{Q}\right) \frac{G M_*^2 R_p^5}{a^6} \frac{n}{\Omega_{\mathrm{rot}}} \sin\varepsilon
\end{align}
where $n = \sqrt{G M_* / a^3}$ is the mean motion and $C_p = I_p / (M_p R_p^2) \approx 0.25$ is the moment-of-inertia coefficient.

\subsection{Observational Selection Function}
\begin{equation}
W_{\mathrm{transit}} = P_{\mathrm{geo}} \cdot P_{\mathrm{det}} = \left(\frac{R_* + R_p}{a}\right) \cdot \frac{1}{1 + \exp\left(-\frac{\mathrm{SNR} - 7.1}{1.5}\right)}
\end{equation}
